%
% Auteur : Olivier Cots
% Date   : 17 avril 2019
%
% styles (environnement) latex
%

% Version noir et blanc, essayer pourquoi pas : 
% https://tex.stackexchange.com/questions/385948/how-to-set-theorem-title-in-middle-of-box

%---------------------------------------------------------------------------------------------------------------
% Environnement todolist, que l'on peut exclure dans le main. Il permet de noter des choses a faire
%
\newenvironment{todolist}{}{}

\ifSLIDES

    %---------------------------------------------------------------------------------------------------------------
    %---------------------------------------------------------------------------------------------------------------
    % Beamer
    
    % gestion du fait qu'il n'y a pas de chapitre :
    %\newcounter{chapter}
    %\regtotcounter{chapter}
    \newcounter{chapter}
    \makeatletter
    \@addtoreset{section}{chapter}
    %\@addtoreset{subsection}{section}
    %\@addtoreset{subsubsection}{subsection}
    \makeatother
    
    \renewcommand \thechapter {\arabic{chapter}}
    \renewcommand \thesection {\thechapter.\arabic{section}}
    %\renewcommand\thesubsection   {\thesection.\@arabic\c@subsection}
    %\renewcommand\thesubsubsection{\thesubsection .\@arabic\c@subsubsection}

    %---------------------------------------------------------------------------------------------------------------
    % Beamer template
    %\usefonttheme{serif}
    %\usetheme{PaloAlto}
    %\usecolortheme[named=ForestGreen]{structure}
    \setbeamertemplate{theorems}[numbered]
    %\setbeamertemplate{theorems}[ams style]
    \setbeamertemplate{navigation symbols}{}
    %\usefoottemplate{\hfill\insertframenumber\,/\,\inserttotalframenumber}
    %\setbeamertemplate{footline}[page number]
    
    \setbeamertemplate{section in toc}[sections numbered]
    \setbeamertemplate{section in toc}{\thechapter.\inserttocsectionnumber.~\inserttocsection}
    \setbeamertemplate{subsection in toc}[subsections numbered]
    \setbeamertemplate{subsection in toc}{\hspace{1.9em}\thechapter.\inserttocsectionnumber.\inserttocsubsectionnumber.~\inserttocsubsection\par}
        
    \setbeamercolor{structure}{fg=black} % titres en noir

    \addtobeamertemplate{frametitle}{}{\vspace*{-5pt}}

    \setbeamertemplate{caption}[numbered] % On numérote les légendes des figures

    %---------------------------------------------------------------------------------------------------------------
    % bibliographie
    %\setbeamertemplate{bibliography item}{\insertbiblabel} % pour avoir des numéros
    
    %---------------------------------------------------------------------------------------------------------------
    % size of margin
    \setbeamersize{text margin left=3mm,text margin right=3mm} 

    %---------------------------------------------------------------------------------------------------------------
    % Couleurs
    \ifBANDW % black and white

        \definecolor{blue}{rgb}{0,0,0}
        \definecolor{red}{rgb}{0,0,0}
        \definecolor{green}{rgb}{0,0,0}
        \definecolor{greenPreuve}{rgb}{0,0,0}
        \definecolor{grey}{rgb}{0.85,0.85,0.85}
        \definecolor{black}{rgb}{0,0,0}
    
        %\definecolor{linkcol}{RGB}{242,157,44} % orange
        \definecolor{linkcol}{RGB}{0,0,0} % orange plus sombre
        %\definecolor{linkcol}{rgb}{0.06,0.06,0.6} % blue
        \definecolor{citecol}{rgb}{0,0,0} % red
        
        %---------------------------------------------------------------------------------------------------------------
        % Couleurs des traits de separation du titre avec le contenu dans les slides
        \definecolor{cfondzero}{rgb}{0,0,0} %{0.0,0.2,0.6}
        \definecolor{cfondun}{rgb}{0,0,0}
        \definecolor{cfonddeux}{rgb}{0,0,0}
        \definecolor{cfondtrois}{rgb}{0,0,0}
        \definecolor{cfondquatre}{rgb}{0,0,0}
        \definecolor{cfondblanc}{rgb}{1.0,1.0,1.0}  
        \definecolor{cfondtitre}{rgb}{0.95,0.95,0.95}
            
        % Couleurs des marqueurs de debut et fin des preuves
        \newcommand{\colorPreuve}{black}
        
        \newtcolorbox{myboxtitle}{colback=cfondtitre,colframe=cfondtitre!80!black,coltext=black}
        
    \else
    
        \ifLIGHT
        
            \definecolor{blue}{rgb}{0.06,0.06,0.6}
            \definecolor{red}{rgb}{0.6,0.04,0.04}
            \definecolor{green}{rgb}{0.06,0.6,0.06}
            \definecolor{greenPreuve}{rgb}{0.0,0.3,0.0}
            \definecolor{grey}{rgb}{0.85,0.85,0.85}
            \definecolor{black}{rgb}{0,0,0}
        
            %\definecolor{linkcol}{RGB}{242,157,44} % orange
            \definecolor{linkcol}{RGB}{213,138,39} % orange plus sombre
            %\definecolor{linkcol}{rgb}{0.06,0.06,0.6} % blue
            \definecolor{citecol}{rgb}{0.6,0.04,0.04} % red
            
            %---------------------------------------------------------------------------------------------------------------
            % Couleurs des traits de separation du titre avec le contenu dans les slides
            \definecolor{cfondzero}{rgb}{0.2,0.6,0.6} %{0.0,0.2,0.6}
            \definecolor{cfondun}{rgb}{0.7,0.4,0.0}
            \definecolor{cfonddeux}{rgb}{0.5,0.5,0.0}
            \definecolor{cfondtrois}{rgb}{0.6,0,0}
            \definecolor{cfondquatre}{rgb}{0.0,0.2,0.6}
            \definecolor{cfondblanc}{rgb}{1.0,1.0,1.0}
            \definecolor{cfondtitre}{rgb}{0.2,0.4,0.6}  
            
            \newtcolorbox{myboxtitle}{colback=cfondtitre!10!white,colframe=cfondtitre!60!white,coltext=black}
        
        \else
        
            % version Retro N7
            \definecolor{blue}{rgb}{0.1,0.1,0.6}
            \definecolor{red}{rgb}{0.6,0.1,0.1}
            \definecolor{green}{rgb}{0.0,0.4,0.0}
            \definecolor{greenPreuve}{rgb}{0.0,0.3,0.0}
            \definecolor{grey}{rgb}{0.85,0.85,0.85}
            \definecolor{black}{rgb}{0,0,0}
            
            %\definecolor{linkcol}{RGB}{242,157,44} % orange
            \definecolor{linkcol}{RGB}{213,138,39} % orange plus sombre
            %\definecolor{linkcol}{rgb}{0.06,0.06,0.6} % blue
            \definecolor{citecol}{rgb}{0.6,0.04,0.04} % red
            
            %---------------------------------------------------------------------------------------------------------------
            % Couleurs des traits de separation du titre avec le contenu dans les slides
            % Version retro N7
            \definecolor{cfondzero}{rgb}{0.2,0.6,0.6} %{0.0,0.2,0.6}
            \definecolor{cfondun}{rgb}{0.5,0.3,0.0}
            \definecolor{cfonddeux}{rgb}{0.3,0.3,0.0}
            \definecolor{cfondtrois}{rgb}{0.5,0,0}
            \definecolor{cfondquatre}{rgb}{0.0,0.2,0.4}
            \definecolor{cfondblanc}{rgb}{1.0,1.0,1.0}
            \definecolor{cfondtitre}{rgb}{0.2,0.4,0.5}
    
            \newtcolorbox{myboxtitle}{colback=cfondtitre,colframe=black,coltext=white}
            
        \fi % light or dark
        
        % Couleurs des marqueurs de debut et fin des preuves
        \newcommand{\colorPreuve}{greenPreuve}
        
    \fi % end Black and white or colored

    %------------------------------------------------------------------------------------------------------------%
    %Page counter
    \newcommand*{\myfont}[1][ptm]{\fontfamily{#1}\selectfont}
    %\DeclareTextFontCommand{\textmyfont}{\myfont[ptm]}
    
    %\usepackage{totcount}
    %\newtotcounter{pages}
    %\setcounter{pages}{1}
    \newcommand*{\compteurPage}[1][red]{
        \begin{tikzpicture}[inner sep=0pt, baseline=(base)]
            \def\epaisseurTrait{0.15em};
            \def\rayonExterne{0.55em};
            \coordinate (O) at (0em,0em);
            \coordinate[shift={(0em,\rayonExterne)}] (A) at (O);
            \pgfmathsetmacro{\myLastPage}{\insertdocumentendpage};
            \pgfmathsetmacro{\myCurrentPage}{\insertpagenumber};
            \pgfmathsetmacro{\myCurrentFrame}{\insertframenumber};
            %\pgfmathsetmacro{\myLastPage}{\inserttotalframenumber};
            %\pgfmathsetmacro{\myCurrentPage}{\insertframenumber};
            \draw[black] (O) node { \scriptsize \bf \myfont{\textcolor{#1}{\myCurrentFrame}}};
            \pgfmathsetmacro{\progress}{(\myCurrentPage/\myLastPage)};
            \draw[grey,line width=\epaisseurTrait,rotate=90] (A) arc (0:360:\rayonExterne);
%            \draw[#1  ,line width=\epaisseurTrait,rotate=90] (A) arc (0:-360*\progress:\rayonExterne);
            \node (base) at (0,-.25ex) {};
        \end{tikzpicture}
    }

    %---------------------------------------------------------------------------------------------------------------
    % Les itemize
    \setlist{itemsep=0.25em}
    \setitemize[1]{label={\small \textbullet}}
    \setitemize[2]{label={\small -}}

\else

    %---------------------------------------------------------------------------------------------------------------
    %---------------------------------------------------------------------------------------------------------------
    % Poly

    %---------------------------------------------------------------------------------------------------------------
    % Couleurs
    \definecolor{blue}{rgb}{0.06,0.06,0.6}
    \definecolor{red}{rgb}{0.6,0.04,0.04}
    \definecolor{green}{rgb}{0.06,0.6,0.06}
    \definecolor{grey}{rgb}{0.85,0.85,0.85}
    \definecolor{black}{rgb}{0,0,0}

    \definecolor{linkcol}{rgb}{0.06,0.06,0.6} % blue
    \definecolor{citecol}{rgb}{0.6,0.04,0.04} % red
    
    % Couleurs Daniel
    \definecolor{ForestGreen}   {cmyk}{0.91,0,0.88,0.12}
    \definecolor{MidnightBlue}  {cmyk}{0.98,0.13,0,0.43}
    
    % Couleurs des marqueurs de debut et fin des preuves
    \newcommand{\colorPreuve}{black}
   
    % Couleurs des titres
    \newcommand{\colorSection}{black}
    \newcommand{\colorSubSection}{black}
    \newcommand{\colorSubSubSection}{black}
    \newcommand{\colorChapter}{black}
    
    %---------------------------------------------------------------------------------------------------------------
    % surcharge de section
    \ifFICHE  

%        \tcbset{
%            orange/.style={
%                enhanced,
%                colback=orange!10,
%                colframe=orange,
%                fonttitle=\bfseries,
%                colbacktitle=orange!40,
%                coltitle=black,
%                attach boxed title to top left={%
%                    yshift*=-\tcboxedtitleheight/2, 
%                    xshift=5mm},
%                boxed title style={colframe=gray}
%            }
%        }
          
%        \titleformat{\chapter}[display]
%        {\normalfont\Large\bfseries}
%        {}
%        {0em}
%        {%
%        \begin{tcolorbox}[orange, halign=center, valign=center, title=test, overlay={\node[draw=gray, fill=orange!40, rounded corners, anchor=east, font=\small\bfseries] at ([xshift=-5mm]frame.south east){2016-2017};}]%, title=Title]
%        \chaptername\space \thechapter #1
%        \end{tcolorbox}%
%        }
  
        % ----------------
        % Section
        \newcommand{\setcolorsectionheading}[2]{%
            %\tikzstyle{sectionstyle}=[draw=#1, text=#2, rounded corners, fill=#1!10, ultra thick, anchor=west, outer sep=0pt]
            \tikzset{sectionstyle_node/.append style={draw=#1, text=#2, fill=#1!10}}
            \tikzset{sectionstyle_draw/.append style={color=#1}}
        }
        
        \tikzstyle{sectionstyle_node}=[draw=orange, text=black, rounded corners=0.0cm, fill=orange!5, ultra thick,
        anchor=west, outer sep=0pt, minimum width=\textwidth, text width=0.96\textwidth, minimum height=1.8cm]
        \tikzstyle{sectionstyle_draw}=[color=orange, ultra thick]
  
        \newcommand{\sectionbreak}{\clearpage} % On saute une page à chaque nouvelle section
        
        \titleformat{\section}
        {\normalfont\Large\bfseries}{}{0em}
        {%
        %\setbox0=\hbox{\large\bfseries\thesection}%
        \begin{tikzpicture}
            %\draw[sectionstyle_draw] (0,0) -- ++(0:\linewidth);
            \node[sectionstyle_node] {\thesection\hspace{.5em}\hangindent\wd0\strut#1\strut};
        \end{tikzpicture}%
        } 
        
        \titleformat{name=\section,numberless}[display]
        {\normalfont\Large\bfseries}{}{0em}
        {%
        \begin{tikzpicture}
            \node[sectionstyle_node] {\strut#1\strut};
        \end{tikzpicture}%
        }

        %\titlespacing*{\section}{0pt}{4.5ex plus 1ex minus .2ex}{2ex plus .2ex}
        \titlespacing*{\section}{0pt}{*4}{*3}
        
        % ----------------
        % Subsection
        \titlespacing*{\subsection}{0pt}{*3}{*2}
        \renewcommand{\colorSubSection}{MidnightBlue}
        %\titleformat{\subsection}{\normalfont\large\bfseries}{{\color{MidnightBlue}\thesubsection\hspace{.5em}\hangindent\wd0\strut#1\strut}}{0em}{} 
        
        % ----------------
        % Subsubsection
        \titlespacing*{\subsubsection}{0pt}{*3}{*2}
        \renewcommand{\colorSubSubSection}{ForestGreen}
        
    \fi
    
    %---------------------------------------------------------------------------------------------------------------
    % Pour le style du manuscrit
    % Clear Header Style on the Last Empty Odd pages
    \makeatletter
    \def\cleardoublepage{
        \clearpage
        \if@twoside 
            \ifodd
                \c@page
            \else%
                \hbox{}%
                \thispagestyle{empty}%              % Empty header styles
                \newpage%
                \if@twocolumn\hbox{}\newpage
                \fi
            \fi
        \fi}
    \makeatother
    
    % minitoc
    \setcounter{minitocdepth}{2}
    \setlength{\mtcindent}{18pt} 
    \renewcommand{\mtcfont}{\small\rm} 
    \renewcommand{\mtcSfont}{\small\bf}
    
    % table des matieres et compteurs de titres
    \setcounter{secnumdepth}{4}
    \setcounter{tocdepth}{2} % ou 1
    
    % centered page environment
    \newenvironment{vcenterpage}
    {\newpage\vspace*{\fill}\thispagestyle{empty}\renewcommand{\headrulewidth}{0pt}}
    {\vspace*{\fill}}
    
    % nicer backref links
    \renewcommand*{\backref}[1]{}
    \renewcommand*{\backrefalt}[4]{%
    \ifcase #1 %
    %
    \or
    $\hookleftarrow$~#2.%
    \else
    $\hookleftarrow$~#2.%
    \fi}
    \renewcommand*{\backrefsep}{, }
    \renewcommand*{\backreftwosep}{ et~}
    \renewcommand*{\backreflastsep}{ et~}
    
    %---------------------------------------------------------------------------------------------------------------
    % Les itemize
    \frenchbsetup{CompactItemize=false}
    \setlist{itemsep=0.0em}
    \newcommand{\myitem}{{\raisebox{0pt}{\textbullet}}}
    \setitemize[1]{label={\small \protect\myitem}}
    \setitemize[2]{label=-}
                                   
\fi

%---------------------------------------------------------------------------------------------------------------
% Configuration du pdf
\hypersetup
{
bookmarksopen=true,
pdfauthor="Olivier COTS",           % auteur du document
pdfstartview={XYZ null null 1},
pdfhighlight=/O,                    % effet d'un clic sur un lien hypertexte
colorlinks=true,                    % couleurs sur les liens hypertextes
linkcolor=linkcol,                  % couleur des liens hypertextes internes
citecolor=citecol,                  % couleur des liens pour les citations
urlcolor=linkcol                    % couleur des liens pour les url
}

%---------------------------------------------------------------------------------------------------------------
% Compteur pour les problèmes
\newcounter{problem}%[chapter]
\regtotcounter{problem}
\newcommand{\tagProblem}[1][]{%
    \addtocounter{problem}{1}
%    \tag{P\theproblem#1}
    \tag{P$_\theproblem$#1}
}

%---------------------------------------------------------------------------------------------------------------
% Les itemize : OLD ?
%
% \newcommand*\circled[1]{\tikz[baseline=(char.base)]{
%     \node[shape=circle,draw,inner sep=2pt] (char) {#1};}}

% \newcommand{\pgftextcircled}[1]{
%     \setbox0=\hbox{#1}%
%     \dimen0\wd0%
%     \divide\dimen0 by 2%
%     \begin{tikzpicture}[baseline=(a.base)]%
%         \useasboundingbox (-\the\dimen0,0pt) rectangle (\the\dimen0,1pt);
%         \node[circle,draw,outer sep=0pt,inner sep=0.1ex] (a) {#1};
%     \end{tikzpicture}
% }

%---------------------------------------------------------------------------------------------------------------
% Listings
\definecolor{mymauve}{rgb}{0.58,0,0.82}
\lstset{ %
  backgroundcolor=\color{white},   % choose the background color; you must add \usepackage{color} or \usepackage{xcolor}
  basicstyle=\small,               % the size of the fonts that are used for the code
  breakatwhitespace=false,         % sets if automatic breaks should only happen at whitespace
  breaklines=true,                 % sets automatic line breaking
  captionpos=b,                    % sets the caption-position to bottom
  commentstyle=\color{blue},       % comment style
  deletekeywords={...},            % if you want to delete keywords from the given language
  escapeinside={\%*}{*)},          % if you want to add LaTeX within your code
  extendedchars=true,              % lets you use non-ASCII characters; for 8-bits encodings only, does not work with UTF-8
  frame=single,                    % adds a frame around the code
  keepspaces=true,                 % keeps spaces in text, useful for keeping indentation of code (possibly needs columns=flexible)
  keywordstyle=\color{red},        % keyword style
  otherkeywords={*,...},           % if you want to add more keywords to the set
  numbers=none,                    % where to put the line-numbers; possible values are (none, left, right)
  numbersep=5pt,                   % how far the line-numbers are from the code
  numberstyle=\tiny\color{gray},   % the style that is used for the line-numbers
  rulecolor=\color{black},         % if not set, the frame-color may be changed on line-breaks within not-black text (e.g. comments (green here))
  showspaces=false,                % show spaces everywhere adding particular underscores; it overrides 'showstringspaces'
  showstringspaces=false,          % underline spaces within strings only
  showtabs=false,                  % show tabs within strings adding particular underscores
  stepnumber=2,                    % the step between two line-numbers. If it's 1, each line will be numbered
  stringstyle=\color{mymauve},     % string literal style
  tabsize=2,                       % sets default tabsize to 2 spaces
  %title=\lstname,                 % show the filename of files included with \lstinputlisting; also try caption instead of title
  frameround=tttt                  % frameround=⟨t|f⟩⟨t|f⟩⟨t|f⟩⟨t|f⟩ ffff
                                   % The four letters designate the top right, bottom right, bottom left and top left corner. 
                                   % In this order. t makes the according corner round.
}

%---------------------------------------------------------------------------------------------------------------
%---------------------------------------------------------------------------------------------------------------
% Environnements : theorem, proposition, etc

\ifSLIDES
    \def\bskip{10pt}
    \def\askip{10pt}
\else
    \def\bskip{10pt}
    \def\askip{10pt}
\fi

\ifSLIDES

    %---------------------------------------------------------------------------------------------------------------
    %---------------------------------------------------------------------------------------------------------------
    % Beamer
    
    % Beamer: to have text on the same line as the title of the environment
    \makeatletter
    \setbeamertemplate{theorem begin}{%
    %  \begin{\inserttheoremblockenv}% removed
        {%
        \ifx\inserttheoremnumber\@empty%
            \textcolor{black}{\bfseries\inserttheoremname}%
        \else
            \textcolor{black}{\bfseries\inserttheoremname\inserttheoremnumber}%
        \fi%
        \ifx\inserttheoremaddition\@empty%
        \else 
            \space(\inserttheoremaddition)%
        \fi%
        \inserttheorempunctuation~%
        % new
      }%
    }
    
    \setbeamertemplate{theorem end}{}
    \makeatother

    %---------------------------------------------------------------------------------------------------------------
    % Couleurs
    %
    % Version retro N7
    \ifBANDW % black and white
        \newcommand\colbackdef{white}\newcommand\colframedef{grey!60!white}\newcommand\coltitledef{black}% Definition
        \newcommand\colbackthm{white}\newcommand\colframethm{grey!40!black}\newcommand\coltitlethm{white} % Theorem
        \newcommand\colbackpro{white}\newcommand\colframepro{grey!65!black}\newcommand\coltitlepro{white}% Proposition
        \newcommand\colbackcor{white}\newcommand\colframecor{grey!65!black}\newcommand\coltitlecor{white}% Corollaire   
    \else
        \ifLIGHT
            \newcommand\colbackdef{white}\newcommand\colframedef{blue!20!white}\newcommand\coltitledef{black}% Definition
            \newcommand\colbackthm{white}\newcommand\colframethm{red!20!white}\newcommand\coltitlethm{black} % Theorem
            \newcommand\colbackpro{white}\newcommand\colframepro{green!20!white}\newcommand\coltitlepro{black}% Proposition
            \newcommand\colbackcor{white}\newcommand\colframecor{green!20!white}\newcommand\coltitlecor{black}% Corollaire
        \else
            \newcommand\colbackdef{white}\newcommand\colframedef{blue!90!white} \newcommand\coltitledef{white} % Definition
            \newcommand\colbackthm{white}\newcommand\colframethm{red!90!white}  \newcommand\coltitlethm{white} % Theorem
            \newcommand\colbackpro{white}\newcommand\colframepro{green!90!white}\newcommand\coltitlepro{white} % Proposition
            \newcommand\colbackcor{white}\newcommand\colframecor{green!90!white}\newcommand\coltitlecor{white} % Corollaire
        \fi
    \fi
    
    %---------------------------------------------------------------------------------------------------------------
    % Style
    
    % encadre formules
    \tcbset{
        every box/.style={highlight math style={enhanced, colback=yellow!25!white,arc=4pt, drop fuzzy shadow, size=minimal, boxsep=5pt}},
        defstyle/.style={enhanced jigsaw, breakable, fonttitle=\bfseries\upshape, fontupper=\slshape, colback=\colbackdef, colframe=\colframedef, separator sign dash,
                       before skip=\bskip,after skip=\askip, fontupper=, coltitle=\coltitledef},
        theostyle/.style={enhanced jigsaw, breakable, fonttitle=\bfseries\upshape, fontupper=\slshape, colback=\colbackthm, colframe=\colframethm, separator sign dash,
                       before skip=\bskip,after skip=\askip, coltitle=\coltitlethm},
        propstyle/.style={enhanced jigsaw, breakable, fonttitle=\bfseries\upshape, fontupper=\slshape, colback=\colbackpro, colframe=\colframepro, separator sign dash,
                       before skip=\bskip,after skip=\askip, coltitle=\coltitlepro},
        corstyle/.style={enhanced jigsaw, breakable, fonttitle=\bfseries\upshape, fontupper=\slshape, colback=\colbackcor, colframe=\colframecor, separator sign dash,
                       before skip=\bskip,after skip=\askip, coltitle=\coltitlecor},
        rmkstyle/.style={enhanced jigsaw, breakable, blanker, left=5mm, before skip=\bskip, after skip=\askip, borderline west={1mm}{0pt}{red}},
        hypstyle/.style={enhanced jigsaw, breakable, blanker, left=5mm, before skip=\bskip, after skip=\askip, borderline west={1mm}{0pt}{black}},
        difstyle/.style={enhanced jigsaw, breakable, blanker, left=5mm, before skip=\bskip, after skip=\askip},
    }


\else

    %---------------------------------------------------------------------------------------------------------------
    %---------------------------------------------------------------------------------------------------------------
    % Poly

    %---------------------------------------------------------------------------------------------------------------
    % Couleurs
    %
    % Version poly plus sobre
    \newcommand\colbackdef{white}\newcommand\colframedef{blue!20!white}\newcommand\coltitledef{black}% Definition
    \newcommand\colbackthm{white}\newcommand\colframethm{red!20!white}\newcommand\coltitlethm{black} % Theorem
    \newcommand\colbackpro{white}\newcommand\colframepro{green!20!white}\newcommand\coltitlepro{black}% Proposition
    \newcommand\colbackcor{white}\newcommand\colframecor{green!20!white}\newcommand\coltitlecor{black}% Corollaire
    
    % Version poly avec plus de couleurs
    %\newcommand\colbackdef{blue!5!white} \newcommand\colframedef{blue!75!black}  \newcommand\coltitledef{white} % Definition
    %\newcommand\colbackthm{red!10!white}   \newcommand\colframethm{red!75!black}   \newcommand\coltitlethm{white} % Theorem
    %\newcommand\colbackpro{green!10!white} \newcommand\colframepro{green!75!black} \newcommand\coltitlepro{white} % Proposition
    %\newcommand\colbackcor{green!10!white} \newcommand\colframecor{green!75!black} \newcommand\coltitlecor{white} % Corollaire

    %---------------------------------------------------------------------------------------------------------------
    % Style
    
    % encadre formules
    \tcbset{
        every box/.style={highlight math style={enhanced, colback=yellow!25!white,arc=4pt, drop fuzzy shadow, size=minimal, boxsep=5pt}},
        defstyle/.style={enhanced jigsaw, breakable, fonttitle=\bfseries\upshape, fontupper=\slshape, colback=\colbackdef, colframe=\colframedef, separator sign dash,
                       before skip=\bskip,after skip=\askip, fontupper=, coltitle=\coltitledef},
        theostyle/.style={enhanced jigsaw, breakable, fonttitle=\bfseries\upshape, fontupper=\slshape, colback=\colbackthm, colframe=\colframethm, separator sign dash,
                       before skip=\bskip,after skip=\askip, coltitle=\coltitlethm},
        propstyle/.style={enhanced jigsaw, breakable, fonttitle=\bfseries\upshape, fontupper=\slshape, colback=\colbackpro, colframe=\colframepro, separator sign dash,
                       before skip=\bskip,after skip=\askip, coltitle=\coltitlepro},
        corstyle/.style={enhanced jigsaw, breakable, fonttitle=\bfseries\upshape, fontupper=\slshape, colback=\colbackcor, colframe=\colframecor, separator sign dash,
                       before skip=\bskip,after skip=\askip, coltitle=\coltitlecor},
        rmkstyle/.style={enhanced jigsaw, breakable, blanker, left=5mm, before skip=\bskip, after skip=\askip, borderline west={1mm}{0pt}{red}},
        hypstyle/.style={enhanced jigsaw, breakable, blanker, left=5mm, before skip=\bskip, after skip=\askip, borderline west={1mm}{0pt}{black}},
        difstyle/.style={enhanced jigsaw, breakable, blanker, left=5mm, before skip=\bskip, after skip=\askip},
    }

\fi

% ----------------------
% Avec tcolorbox

\newtcbtheorem[number within=section]{mytheorem}{Théorème}{theostyle}{thm}
\newtcbtheorem[use counter from=mytheorem]{mydefinition}{Définition}{defstyle}{def}
\newtcbtheorem[use counter from=mytheorem]{myproposition}{Proposition}{propstyle}{prop}
\newtcbtheorem[use counter from=mytheorem]{mycorollary}{Corollaire}{corstyle}{cor}

\ifSLIDES
    \resetcounteronoverlays{tcb@cnt@mytheorem} % to avoid theorem numbers from changing within a frame
    \resetcounteronoverlays{tcb@cnt@mydefinition}
    \resetcounteronoverlays{tcb@cnt@myproposition}
    \resetcounteronoverlays{tcb@cnt@mycorollary}
\fi

% ----------------------
% Avec newtheorem
%\ifSLIDES
%    \newtheorem{mylemma}{Lemme}[]    % Lemme
%\else
    \newtheorem{mylemma}{Lemme}[section]    % Lemme
%\fi
\newtheorem*{mylemma*}{Lemme}    % Lemme

%---------------------------------------------------------------------------------------------------------------
% Environnement : Proof
%
\newcommand{\newstep}{\textasteriskcentered}           % Pour marquer une nouvelle etape dans la preuve

\newcommand*{\QEDA}{\hfill\ensuremath{\blacksquare}}%
\newcommand*{\QEDB}{\hfill\ensuremath{\square}}%

\makeatletter
\renewenvironment{proof}[1][{\color{\colorPreuve} $\RHD$}]{
\par\pushQED{{\color{\colorPreuve}\QEDA}}\normalsize\normalfont\topsep6\p@\@plus6\p@\relax\trivlist\item[\hskip\labelsep\itshape#1\@addpunct{}]\ignorespaces
}{
    \popQED\endtrivlist\@endpefalse
}
\makeatother
%
%\ifSLIDES
%    \tcolorboxenvironment{proof}{blanker, breakable, left=1.5em, before skip=10pt, after skip=10pt, coltext=\colorPreuve}
%\else
    \tcolorboxenvironment{proof}{blanker, breakable, left=1.5em, before skip=\bskip, after skip=\askip, coltext=\colorPreuve}
%\fi

%\ifSLIDES

    %---------------------------------------------------------------------------------------------------------------
    %---------------------------------------------------------------------------------------------------------------
    % Beamer
    
    % proof deb, mil et fin
    %
    % DEB
    \makeatletter
    \newenvironment{proofdeb}[1][{\color{\colorPreuve} $\RHD$}]{
    \par\pushQED{{\color{\colorPreuve}\QEDA}}\normalsize\normalfont\topsep6\p@\@plus6\p@\relax\trivlist\item[\hskip\labelsep\itshape#1\@addpunct{}]\ignorespaces
    }{
        %\popQED
        \endtrivlist\@endpefalse
    }
    \makeatother
    %
    \tcolorboxenvironment{proofdeb}{blanker, breakable, left=1.5em, before skip=\bskip, after skip=\askip, coltext=\colorPreuve}
    
    % MIL
    \makeatletter
    \newenvironment{proofmil}[1][{\color{\colorPreuve} $\RHD$}]{
    \par\pushQED{{\color{\colorPreuve}\QEDA}}\normalsize\normalfont
    }{
        %\popQED
        \endtrivlist\@endpefalse
    }
    \makeatother
    %
    \tcolorboxenvironment{proofmil}{blanker, breakable, left=1.5em, before skip=\bskip, after skip=\askip, coltext=\colorPreuve}
    
    % FIN
    \makeatletter
    \newenvironment{prooffin}[1][]{
    \par\pushQED{{\color{\colorPreuve}\QEDA}}\normalsize\normalfont
    }{
        \popQED\endtrivlist\@endpefalse
    }
    \makeatother
    %
    \tcolorboxenvironment{prooffin}{blanker, breakable, left=1.5em, before skip=\bskip, after skip=\askip, coltext=\colorPreuve}

%\fi

%---------------------------------------------------------------------------------------------------------------
% Environnement : style en cours pour ce qui est defini par un newtheorem
%
% example, myremark, exerciceTHM, corrTHM, 
%

\ifSLIDES

    \newtheoremstyle{mystyleExercice}%          % Name
      {}%                                       % Space above
      {}%                                       % Space below
      {}%                                       % Body font
      {}%                                       % Indent amount
      {}%                                       % Theorem head font
      {.}%                                      % Punctuation after theorem head
      { }%                                      % Space after theorem head, ' ', or \newline
    %  {}%                                      % Theorem head spec (can be left empty, meaning `normal')
      {}%
    
    \theoremstyle{mystyleExercice}
    
\else

    \newtheoremstyle{mystyleExercice}%          % Name
      {}%                                       % Space above
      {}%                                       % Space below
      {}%                                       % Body font
      {}%                                       % Indent amount
      {}%                                       % Theorem head font
      {.}%                                      % Punctuation after theorem head
      { }%                                      % Space after theorem head, ' ', or \newline
    %  {}%                                      % Theorem head spec (can be left empty, meaning `normal')
      {\noindent\textbf{\thmname{#1}\thmnumber{~#2}}\thmnote{\color{black}~(#3)}}%
    
    \theoremstyle{mystyleExercice}

\fi

%---------------------------------------------------------------------------------------------------------------
% Environnement : Exemple
%

\ifSLIDES

    %---------------------------------------------------------------------------------------------------------------
    %---------------------------------------------------------------------------------------------------------------
    % Beamer
    
    \newtheorem{myexample}{Exemple}[section]      % Prend le style en cours
    % \newtheorem{exampleBeamer}{Exemple}[chapter]      % Prend le style en cours
    % \makeatletter
    % \newenvironment{myexample}[1][]
    % {
    %     \begin{exampleBeamer}[#1]
    %     %\par\pushQED{\QEDB}
    % }
    % {
    %     %\popQED\endtrivlist\@endpefalse
    %     \end{exampleBeamer}
    % }
    % \makeatother
    
    \newtheorem*{myexample*}{Exemple}
    % \newtheorem*{exampleBeamer*}{Exemple}
    % \makeatletter
    % \newenvironment{myexample*}[1][]
    % {
    %     \begin{exampleBeamer*}[#1]
    %     %\par\pushQED{\QEDB}
    % }
    % {
    %     %\popQED\endtrivlist\@endpefalse
    %     \end{exampleBeamer*}
    % }
    % \makeatother

\else

    %---------------------------------------------------------------------------------------------------------------
    %---------------------------------------------------------------------------------------------------------------
    % Poly
    
    \newtheorem{example}{Exemple}[section]      % Prend le style en cours
    \makeatletter
    \newenvironment{myexample}[1][]
    {
        \begin{example}[#1]
        \par\pushQED{\QEDB}
    }
    {
        \popQED\endtrivlist\@endpefalse
        \end{example}
    }
    \makeatother

    \newtheorem*{example*}{Exemple}      % Prend le style en cours
    \makeatletter
    \newenvironment{myexample*}[1][]
    {
        \begin{example*}[#1]
        \par\pushQED{\QEDB}
    }
    {
        \popQED\endtrivlist\@endpefalse
        \end{example*}
    }
    \makeatother

\fi

%---------------------------------------------------------------------------------------------------------------
% Environnement : Remarque. C'est un environnement classique encapsule dans un tcolorbox
%

%\ifSLIDES
%    \newtheorem{myremark}{Remarque}
%\else
    \newtheorem{myremark}{Remarque}[section]
%\fi
\newtheorem*{myremark*}{Remarque}
\tcolorboxenvironment{myremark}{rmkstyle}
\tcolorboxenvironment{myremark*}{rmkstyle}

%---------------------------------------------------------------------------------------------------------------
% Environnement : Exercice, Question, Correction, etc
%
\newcommand{\exenotext}{~\\ \vspace{-\baselineskip}}            % Utile avant un enumerate dans un exercice s'il n'y a pas de texte avant
\newcommand{\anoter}{ {\Large \textcolor{orange}{\ding{45}}} }  % La main qui ecrit avant un exercice
\newcommand{\newquestion}{\stepcounter{question}\setcounter{subquestion}{1}}   % Utile pour la correction d'un exo

% Exercice
\newcounter{question}               % Compteur pour les questions
\newcounter{subquestion}           % Compteur pour les sous-questions

\ifSLIDES

    %---------------------------------------------------------------------------------------------------------------
    %---------------------------------------------------------------------------------------------------------------
    % Beamer
    
    %\newtheorem{myexercise}[myexample]{Exercice}
     \newtheorem{exerciceTHM}[myexample]{Exercice}
     \makeatletter
     \newenvironment{myexercise}[1][]
     {
         \begin{exerciceTHM}[#1]
         %\par\pushQED{\QEDB}
         \setcounter{question}{0}
     }
     {
         \setcounter{question}{0}
         %\popQED\endtrivlist\@endpefalse
         \end{exerciceTHM}
     }
    
    %\newtheorem*{myexercise*}{Exercice}
     \newtheorem*{exerciceTHM*}{Exercice}
     \newenvironment{myexercise*}[1][]
     {
         \begin{exerciceTHM*}[#1]
         %\par\pushQED{\QEDB}
         %\setcounter{question}{0}%
     }
     {
         %\setcounter{question}{0}%
         %\popQED\endtrivlist\@endpefalse
         \end{exerciceTHM*}
     }
     \makeatother

\else

    %---------------------------------------------------------------------------------------------------------------
    %---------------------------------------------------------------------------------------------------------------
    % Poly
    
    \newtheorem{exerciceTHM}[example]{{\makebox[0pt][r]{\anoter\hspace*{1ex}}Exercice}}
    \makeatletter
    \newenvironment{myexercise}[1][]
    {
        \begin{exerciceTHM}[#1]
        \par\pushQED{\QEDB}
        \setcounter{question}{0}
    }
    {
        \setcounter{question}{0}
        \popQED\endtrivlist\@endpefalse
        \end{exerciceTHM}
    }
    \makeatother

    \newtheorem*{exerciceTHM*}{{\makebox[0pt][r]{\anoter\hspace*{1ex}}Exercice}}
    \makeatletter
    \newenvironment{myexercise*}[1][]
    {
        \begin{exerciceTHM*}[#1]
        \par\pushQED{\QEDB}
        %\setcounter{question}{0}
    }
    {
        %\setcounter{question}{0}
        \popQED\endtrivlist\@endpefalse
        \end{exerciceTHM*}
    }
    \makeatother

\fi

% Question
\makeatletter
\newenvironment{myquestion}
{
    \normalsize\normalfont\topsep6\p@\@plus6\p@\relax\trivlist
    \item[\stepcounter{question}\bfseries\hspace{0.5em}\hskip\labelsep\thequestion\@addpunct{.}]\ignorespaces
    \setcounter{subquestion}{1}
}
{
}
\makeatother

% Sous Question
\makeatletter
\newenvironment{mysubquestion}
{
    \normalsize\normalfont\topsep6\p@\@plus6\p@\relax\trivlist
    \item[\bfseries\hspace{1.5em}\hskip\labelsep\thequestion\@addpunct{.}\thesubquestion\@addpunct{.}]\ignorespaces
}
{
    \stepcounter{subquestion}
}
\makeatother

% Correction exercice
\newtheorem*{corrTHM}{{\makebox[0pt][r]{$\rhd$\hspace*{1ex}}Correction}}
\makeatletter
\newenvironment{mycorr}[1][]
{
    \begin{corrTHM}[#1]
    \setcounter{question}{0}
    \par\pushQED{\QEDB}
}
{
    \popQED\endtrivlist\@endpefalse
    \setcounter{question}{0}
    \end{corrTHM}
}
\makeatother

%---------------------------------------------------------------------------------------------------------------
% Environnement : Hypothese
\newtheorem{assumption}{Hypothèse}

\newcounter{assumptionx}
\regtotcounter{assumptionx}
\newcommand{\tagAssumption}{%
    H\theassumptionx
}

\newenvironment{myassumption}
{\addtocounter{assumptionx}{1}
    \renewcommand\theassumption{\tagAssumption}
    \addtocounter{assumption}{-1}
    \assumption
}
{
    \endassumption
}

\tcolorboxenvironment{myassumption}{hypstyle}

\ifSLIDES
    \resetcounteronoverlays{assumptionx}
\fi

%---------------------------------------------------------------------------------------------------------------
% Environnement : Question/Difficulte a se poser
%
\newtheorem{questionnement}{}

\newcounter{questionnementx}
\regtotcounter{questionnementx}
\newcommand{\tagquestionnement}{%
Q\thequestionnementx
}

\newenvironment{myquestionnement}
{\addtocounter{questionnementx}{1}
    \renewcommand\thequestionnement{\tagquestionnement}
    \addtocounter{questionnement}{-1}
    \questionnement
}
{
    \endquestionnement
}

\tcolorboxenvironment{myquestionnement}{difstyle}

\ifSLIDES
    \resetcounteronoverlays{questionnementx}
\fi

%---------------------------------------------------------------------------------------------------------------
% Environnement : Question/Difficulte a se poser
%
\newtheorem{difficulty}{}

\newcounter{difficultyx}
\regtotcounter{difficultyx}
\newcommand{\tagDifficulty}{%
D\thedifficultyx
}

\newenvironment{mydifficulty}
{\addtocounter{difficultyx}{1}
    \renewcommand\thedifficulty{\tagDifficulty}
    \addtocounter{difficulty}{-1}
    \difficulty
}
{
    \enddifficulty
}

\tcolorboxenvironment{mydifficulty}{difstyle}

\ifSLIDES
    \resetcounteronoverlays{difficultyx}
\fi


% Exercices colorbox avec record
% \tcbuselibrary{skins,xparse}
%
% to be possible to reference the exercices I make an encapsulation
%\newtheorem{myexercisecb_}{}[section]
%
\newcounter{myexercisecbinner}[section]
\newcounter{myexercisecbouter}[section]
\renewcommand{\themyexercisecbinner}{\thesection.\themyexercisecbouter}
\def\stepexecbcounter{\stepcounter{myexercisecbouter}\refstepcounter{myexercisecbinner}}
%
% 
\NewTColorBox[auto counter, number within=section]{myexercisecb}{+!O{} d<>}%
    {enhanced, 
    breakable, 
    colframe=green!20!white, 
    colback=white, 
    coltitle=green!40!black, 
    fonttitle=\bfseries, 
    lower separated=false,
    title={Exercice~\thetcbcounter~(solution p.~\pageref{mysolution@\thetcbcounter})~:}, 
    label={exercise@\thetcbcounter},
    attach title to upper=~,
    before title={\stepexecbcounter\IfValueT{#2}{\label{ex:#2}}},
    lowerbox=ignored,
    savelowerto=solutions/exercise-\thetcbcounter.tex, 
    record={\string\mysolution{\thetcbcounter}{solutions/exercise-\thetcbcounter.tex}}, #1
}
%
\NewTotalTColorBox{\mysolution}{mm}{%
    enhanced, breakable, colframe=red!20!white, colback=white, coltitle=red!40!black, fonttitle=\bfseries,
    title={Solution de l'Exercice~\ref{exercise@#1} p.~\pageref{exercise@#1}~:},
    phantomlabel={mysolution@#1},
    attach title to upper=~,
}{\input{#2}}
%
\tcbset{no solution/.style={no recording, title={Exercice~\thetcbcounter:}}}

